\section{Introduction}
\label{sec:intro}

This paper answers a question posed by the author of PROJECT spacetimes: can a new gauge
symmetry act on a quantized 16-component Dirac field on the eight-dimensional Minkowski
spacetime of signature $7+1$, with a Lagrangian that is quadratic and quartic in the Dirac
field, such that the effective potential of the field has a minimum at which (i)~a superfluid
of bosonic spin-zero pairs, built from the Dirac field alone, forms; (ii)~the new gauge group
is spontaneously broken; and (iii)~a new electrically neutral scalar boson, which exists only
after the breaking, carries the attraction between the members of the pairs?

The spacetime is $\mathbb R^8$ with the coordinates $X=\{x0,\dots,x7\}$, of which $x4$ is the
time, and the metric $\eta_{1671}=\mathrm{diag}(-1,-1,-1,-1,+1,-1,-1,-1)$. It is contractible,
so it carries a spin structure and Dirac spinors exist; they have $2^{8/2}=16$ components. The
fields of this paper are quantized: their components are complex Grassmann operators.

The answer, worked out below and checked by machine at every step, is yes, with one
qualification. The three requirements are met by the model of Secs.~\ref{sec:pairs}
to~\ref{sec:breaking}. The qualification concerns the word ``gauge'' in ``scalar gauge
boson'': a gauge field is a connection, a one-form, and its quanta are vectors. The two new
neutral bosons that exist only after the breaking are a massive vector, $B_{U2A}$, the analogue
of the Z boson of the electroweak theory, and a scalar, the amplitude mode $\sigma_{U2}$ of
the pair condensate, the analogue of the Higgs boson. The scalar is the one that carries
the pairing interaction, in the precise sense of Sec.~\ref{sec:sigma}.

The electroweak theory~\cite{Weinberg1967} is the template. There the Higgs doublet
$\phi_{\rm SM}$ of hypercharge $Y_{\rm SM}=\tfrac12$ has the covariant derivative, written with
one-forms,
\begin{equation}
\mathrm D^{\rm SM}\phi_{\rm SM}=\big(\mathrm d+\tfrac i2\,g_{\rm SM}\,\boldsymbol\tau_{\rm SM}\cdot\boldsymbol W_{\rm SM}
+i\,g'_{\rm SM}Y_{\rm SM}B_{\rm SM}\big)\phi_{\rm SM},
\label{eq:sm}
\end{equation}
with the Pauli matrices $\boldsymbol\tau_{\rm SM}$, the gauge fields $\boldsymbol W_{\rm SM}$,
$B_{\rm SM}$ and the couplings $g_{\rm SM}$, $g'_{\rm SM}$; its expectation value breaks
$SU(2)\times U(1)$ to the $U(1)$ of electromagnetism, and the bosons W$^\pm$, Z, the photon and
the Higgs boson appear, the Z and the Higgs boson electrically neutral, and both defined only
after the breaking. In this paper the covariant derivative acts on the fermions themselves,
and the field whose expectation value breaks the symmetry is not elementary: it is a pair of
fermions, as in the theory of superconductivity~\cite{BCS1957,Nambu1960} and in the model of
Nambu and Jona-Lasinio~\cite{NJL1961}.

Three facts specific to the $7+1$ spacetime shape the model. First (Sec.~\ref{sec:pairs}),
both charge-conjugation matrices $C^\pm_{1671}$ of this spacetime are symmetric, so a single
Dirac field has no Lorentz-scalar pair; a spin-zero pair needs two fields in the antisymmetric
combination, and the group that preserves that combination is $SU(2)$. Second
(Sec.~\ref{sec:gauge}), the left-handed and the right-handed spinors of $7+1$ are complex
conjugates of each other, so a doublet of one chirality can pair but cannot carry a gauged
$U(1)$ without an anomaly; the gauged doublet is a Dirac doublet. Third (Sec.~\ref{sec:potential}),
the theory is not renormalizable: the quartic couplings have the dimension mass$^{-6}$ and the
gauge couplings mass$^{-2}$, so the momentum cutoff $\Lambda_{U2}$ is part of the model, and we
compute with two cutoffs to show which results depend on it.

The paper is organized as follows. Section~\ref{sec:spacetime} sets out the Dirac matrices,
Sec.~\ref{sec:pairs} the spin-zero pairs, Sec.~\ref{sec:gauge} the new gauge group and its
covariant derivative, Sec.~\ref{sec:lagrangian} the Lagrangian and the self-interacting Dirac
equations, Sec.~\ref{sec:meanfield} their mean-field spinor solutions and the superfluid state,
Sec.~\ref{sec:potential} the effective potential and its minima, Sec.~\ref{sec:breaking} the
breaking of the gauge group and the bosons that it produces, and Sec.~\ref{sec:status} what has
been established and what has not. Appendix~\ref{app:notation} defines every symbol,
Appendix~\ref{app:checks} lists the machine checks, and Appendix~\ref{app:data} the programs and
data. The suffix U2 marks every symbol of the new sector, so that no symbol of this paper can
be mistaken for a symbol of the same letter elsewhere in the project; units are those of the
cutoff, $\Lambda_{U2}=1$, unless a unit is shown.
